Under one protocol, the reimplemented ESN forecasts Lorenz-63 for 9.546 Lyapunov times at 5000 training samples, 1.71 times the tuned delay-coordinate MLP and 2.52 times the GRU, and has the highest mean VPT at every training length tested. On Lorenz-96 ($D=10$) the ESN beats the GRU but is behind the MLP (2.115 versus 2.882, $p=0.018$), so the hypothesis that the ESN wins on both tasks is refuted. Both gaps depend on settings that a comparison must report: the optimiser budget of the MLP moves the Lorenz-63 ratio between 2.74 and 1.69, and on Lorenz-96 a reservoir of 2000 nodes has a higher mean VPT than the tuned MLP (3 seeds, not tested). The ESN is sensitive to spectral radius (best/worst VPT ratios of 4.18 and 10.95) and improves with reservoir size. With 500 training samples it has the highest mean VPT on both tasks, though not significantly above the MLP with 3 seeds. At the main setting the three systems do not differ in the marginal attractor statistics we measure; on Lorenz-96, free runs leave the attractor region without squared features (ESN) or input noise (GRU), in the GRU's case with unchanged VPT. With two systems, 3--5 seeds and small tuning grids, this is a replication-scale comparison, not a ranking of model families.
